Účel a zásady
\begin{itemize}\item Cílem zrychlit tvorbu cvičení, studijních opor a multimédií při zachování akademické kvality a souladu s výukovými cíli; pracujte iterativně: AI navrhne, učitel upraví a validuje, pilot a úpravy.\end{itemize}
Výběr nástrojů — krátký checklist
\begin{itemize}\item Podpora vašich formátů (LMS, QTI/CSV, multimédia), transparentnost modelu a editovatelnost, řízení autorských práv/provenance a datové zabezpečení/nasazení on‑premise.\end{itemize}
Workflow (krok‑za‑krok)
\begin{enumerate}\item Definujte výukové cíle a formát; vytvořte šablonu promptu; generujte první verzi; učitel kontroluje správnost, shodu s cíli, obtížnost, inkluzi a rizika; pilot, sběr metrik a finalizace/integace do LMS.\end{enumerate}
Ukázka a kontrola kvality
\begin{itemize}\item Příklad: učitel nastaví cíle a šablony, AI vygeneruje varianty, náhodná kontrola; zajistit faktickou správnost, respekt k autorským právům, inkluzi, archivaci promptů a měření přínosu (čas, kvalita, spokojenost, výsledky).\end{itemize}